\section{An integrable anticanonical metric on a toroidal base}
\label{g:sec-toroidal-volume}

A smooth metric on the anticanonical bundle of a smooth variety determines
a smooth positive volume form.  Integrating that form along a fibration
produces a density on the base.  The question here is whether this density
is the squared norm of a positively curved anticanonical metric, and
whether its singularities can be corrected without changing divisors on
the original variety.  Adjoint rank one answers the first question.
The comparison between ramification and Jacobian orders answers the
second.

We use additive notation for line bundles.  If a local frame has squared
norm $e^{-\varphi}$, its curvature is $i\partial\bar\partial\varphi$.
For a rational line bundle, choose a positive Cartier multiple and a
nonvanishing local holomorphic frame of that multiple, and divide its
weight by the chosen integer.  Integrability below means
$e^{-\varphi}\in L^1_{\mathrm{loc}}$ for these root weights.  Changing
the holomorphic frame multiplies the squared norm by a locally bounded
positive factor with locally bounded inverse, and therefore does not
change this condition.

A toroidal embedding is a normal variety with an open subset locally
modeled on a toric variety and its dense torus.  Its boundary is the
complement of that open subset.  We use strict toroidal embeddings,
whose irreducible boundary components are normal, and morphisms locally
modeled by toric morphisms.  The full reduced boundary will let us
control discrepancies even when the source variety is singular.

\begin{theorem}[Toroidal volume metric]\label{g:toroidal-volume}
Let $X$ and $Y$ be smooth connected projective complex varieties, with
$n=\dim X$ and $d=\dim Y>0$, and let $L=-K_X$ carry a smooth Hermitian metric $h$ of
semipositive curvature $\alpha$.  Suppose
\[
 p:W\longrightarrow X,\qquad f:W\longrightarrow Y
\]
are projective morphisms, where $W$ is normal and integral, $p$ is
birational, and $f$ is surjective, equidimensional, and toroidal for
strict toroidal embeddings.  Choose a smooth projective resolution
$\rho:Z\to W$, and put $\pi=p\rho$, $g=f\rho$, and $E_Z=K_{Z/X}$.
Assume
\begin{equation}\label{g:volume-rank-zero}
 \operatorname{rk}g_*\mathcal O_Z(E_Z)=1.
\end{equation}
For each prime divisor $Q\subset Y$ and each prime $E\subset W$
dominating $Q$, set
\[
 a_E=\operatorname{ord}_E K_{W/X},\qquad
 e_E=\operatorname{ord}_E f^*Q,
\]
where the Jacobian divisor $K_{W/X}=K_W-p^*K_X$ is defined in
codimension one.  Define
\begin{equation}\label{g:volume-correction}
 a_Q=\max\left\{0,\min_{f(E)=Q}
             \left(\frac{1+a_E}{e_E}-1\right)\right\},
 \qquad A_0=\sum_Q a_Q Q.
\end{equation}
Then $A_0$ is an effective rational divisor with finite support, and
$f^*A_0$ is $p$-exceptional.  The rational line bundle $-K_Y+A_0$
has a singular Hermitian metric of semipositive curvature whose squared
frame norms are locally integrable without a divisor weight.  This
metric is smooth on a nonempty Zariski open subset of $Y$.

Moreover, let $U\subset Y$ be a nonempty Zariski open set where $g$
is smooth and formation of the adjoint direct image
$g_*(K_{Z/Y}+\pi^*L)$ commutes with base change.  If there are a
nonempty ordinary open patch
$\mathcal V\subset g^{-1}(U)\setminus\operatorname{Supp}E_Z$,
a K\"ahler form $\omega_Y$, and a constant $c>0$ such that
\begin{equation}\label{g:volume-horizontal-patch}
 \pi^*\alpha\geq c\,g^*\omega_Y\qquad\text{on }\mathcal V,
\end{equation}
then the constructed metric is strictly positively curved on a
nonempty ordinary open subset of $Y$.
\end{theorem}

The rank in \eqref{g:volume-rank-zero} is the full generic adjoint rank
at exponent zero.  It identifies the entire direct image with a line
on a suitable open set.  In particular, the theorem does not obtain
positivity by restricting a direct-image metric to an arbitrary
subline.  Likewise, \eqref{g:volume-horizontal-patch} is a hypothesis
on one patch; its conclusion is strict curvature on an ordinary open,
not a global K\"ahler lower bound.

The theorem assumes a toroidal equidimensional model rather than
constructing one.  Birational weak toroidalization and projective
subdivisions provide such models in the applications; these are the
constructions of Abramovich--Denef--Karu and Abramovich--Karu
\cite{ADK2013,AK2000}.  No cover making the fibers reduced is part of
the theorem.  We retain the multiplicities $e_E$ in the correction.

\subsection{The correction and its discrepancies}

We first isolate the algebraic reason that the correction is disposable
on $X$ and integrable on every higher model.  The transverse map $y_1=u_1^{e_E}$ sends a density of order
$|u_1|^{2a_E}$ to one of order
$|y_1|^{2(a_E-e_E+1)/e_E}$. This explains the coefficient in
\eqref{g:volume-correction}. After the correction, on a normal-crossing
resolution $Z^\prime\to W$, with composite $g^\prime:Z^\prime\to Y$,
the upstairs density will have factors
$|z_j|^{2\gamma_j}$ with
$\gamma_j=\operatorname{coeff}_{\{z_j=0\}}
(K_{Z'/X}-(g')^*A_0)$. Such factors are integrable exactly when
$\gamma_j>-1$. New exceptional divisors can be centered in boundary
intersections, so checking only the primes already on $W$ is insufficient.
The following lemma controls all the new exponents by comparing with the
full toroidal boundary. The canonical divisors in
this discussion are chosen compatibly with a rational top form on
$X$; all Jacobian differences are consequently unambiguous.

For a rational Weil divisor $\Theta$ with $K_W+\Theta$ rationally
Cartier, its crepant pullback to a smooth model $\tau:W'\to W$ is
defined by $K_{W'}+\Theta_{W'}=\tau^*(K_W+\Theta)$.
The pair is \emph{sub-klt} if every coefficient of $\Theta_{W'}$
is less than one on every such model; negative coefficients are
allowed.  Equivalently, every divisorial log discrepancy
$A_{W,\Theta}(\operatorname{ord}_F)=1-\operatorname{coeff}_F\Theta_{W'}$
is positive.  Replacing ``positive'' by ``nonnegative'' defines
sub-log-canonicity, or log canonicity when the boundary is effective.

\begin{lemma}\label{g:volume-subklt}
Let $X,Y$ be smooth connected projective varieties with $\dim Y>0$,
let $W$ be normal and integral, and let $p:W\to X$ be projective
birational and $f:W\to Y$ projective, surjective, equidimensional,
and toroidal for strict toroidal embeddings. Define $A_0$ by
\eqref{g:volume-correction}, and put
\[
 \Theta=-K_{W/X}+f^*A_0.
\]
Then $\Theta\leq0$, the pair $(W,\Theta)$ is sub-klt, and
$f^*A_0$ is $p$-exceptional.  In particular, for every smooth model
$\tau:Z'\to W$,
\begin{equation}\label{g:volume-all-valuations}
 \operatorname{coeff}_F\bigl(K_{Z'/X}-(f\tau)^*A_0\bigr)>-1
\end{equation}
at every prime $F\subset Z'$.
\end{lemma}

This assertion uses only the morphisms and their divisor orders.
Neither a metric nor an adjoint-rank assumption is needed; this will
allow us to use the same discrepancy estimate with invariant direct
images later.

\begin{proof}
At a general point of a prime of $W$, the variety is smooth and the
Jacobian of the morphism to $X$ is holomorphic.  Hence $a_E\geq0$.
At a prime that is not exceptional over $X$, the morphism is an
isomorphism at its general point, so $a_E=0$.

Every vertical prime of $W$ dominates a prime of $Y$.  Indeed, an
equidimensional map of relative dimension $\dim W-\dim Y$ cannot map a
divisor into a set of codimension at least two: its dimension would
then be at most $\dim W-2$.  If $a_Q>0$, the minimum in
\eqref{g:volume-correction} is positive, so every prime over $Q$ has
$a_E>0$ and is $p$-exceptional.  This proves the asserted support of
$f^*A_0$.  It also proves finite support, since there are only finitely
many $p$-exceptional prime divisors.  Rationality follows from the
integer orders in \eqref{g:volume-correction}.

If $E$ dominates $Q$ with $a_Q>0$, its coefficient in $\Theta$
satisfies
\[
 -a_E+e_Ea_Q\leq 1-e_E\leq0.
\]
At a prime above a $Q$ with $a_Q=0$, and at any horizontal prime,
the coefficient is $-a_E\leq0$.  Thus $\Theta\leq0$.

Let $D_W$ be the full reduced toroidal boundary.  The toroidal pair
$(W,D_W)$ is log canonical, $K_W+D_W$ is rationally Cartier, and
$W\setminus D_W$ is smooth.  These assertions can be checked on a
regular subdivision of each toric chart: the logarithmic top form
gives
$K_{\widetilde W}+D_{\widetilde W}=\tau^*(K_W+D_W)$,
and the reduced boundary upstairs has simple normal crossings.
We also have
\[
 K_W+\Theta=p^*K_X+f^*A_0.
\]
Consequently $B=D_W-\Theta$ is an effective rational Cartier divisor
whose support contains $D_W$.  For any divisorial valuation $v$ over
$W$, the log discrepancies satisfy
\begin{equation}\label{g:volume-discrepancy-comparison}
 A_{W,\Theta}(v)=A_{W,D_W}(v)+v(B).
\end{equation}
The full toroidal boundary gives nonnegative log discrepancies;
the effective rational Cartier difference adds a strictly positive term precisely
where that discrepancy could vanish. Both terms on the right are nonnegative.  If the first is zero, the
center of $v$ lies in $D_W$, since the complement is smooth with zero
boundary.  Its center is then contained in the support of the
effective rational Cartier divisor $B$, so $v(B)>0$.  If the first
term is positive there is nothing more to check.  This proves
$A_{W,\Theta}(v)>0$ for every divisorial valuation, which is the
sub-klt assertion.  This argument does not require the support of
$K_{W/X}$ to lie in the toroidal boundary.

Finally the crepant pullback of $\Theta$ to $Z'$ is
\[
 \Theta_{Z'}=(f\tau)^*A_0-K_{Z'/X}.
\]
Its coefficients are less than one by the sub-klt assertion, giving
\eqref{g:volume-all-valuations}.
\end{proof}

\subsection{Positivity of the integrated density}

We next explain why the volume density gives a semipositive metric,
and why the patch in \eqref{g:volume-horizontal-patch} gives strict
curvature.  The analytic input is the smooth direct-image theorem
of Berndtsson \cite[Theorem~1.2]{Berndtsson2009}.  We use its
primitive-representative curvature formula before the additional
product and fiberwise-positivity assumptions considered later in that
paper.

For clarity, consider a proper smooth K\"ahler family
$f:\mathcal Z\to\mathbb D$ over a disc with coordinate $t$,
relative dimension $r>0$, and a smooth semipositive
coefficient line $(\mathcal L,h_{\mathcal L})$.  Assume formation of
$\mathcal E=f_*(K_{\mathcal Z/\mathbb D}+\mathcal L)$ commutes with
base change.  Choose a holomorphic section $u$ of $\mathcal E$
that is Chern-normal at $0$.  Berndtsson's representative construction
\cite[Proposition~4.2]{Berndtsson2009} gives an absolute
$\mathcal L$-valued $(r,0)$-form $\widehat u$ representing $u$, with
\[
 \bar\partial\widehat u=dt\wedge\eta,\qquad
 \partial^{h_{\mathcal L}}\widehat u=dt\wedge\mu,
 \qquad \mu|_{\mathcal Z_0}=0,
 \qquad (\omega\wedge\eta)|_{\mathcal Z_0}=0.
\]
Here $\omega$ is a K\"ahler form on the total space.  Differentiating
the fiber integral with this representative gives
\begin{equation}\label{g:volume-curvature-bound}
 -i\partial\bar\partial\|u\|^2\big|_0
 \ \geq\
 f_*\!\left(i^{r^2}\widehat u\wedge\overline{\widehat u}
       h_{\mathcal L}\wedge i\Theta_{\mathcal L,h_{\mathcal L}}\right)
       \big|_0.
\end{equation}
The omitted term is the nonnegative squared norm of the primitive
$(r-1,1)$-form $\eta$; this is
\cite[formula~(4.8)]{Berndtsson2009}.
The displayed integrand is nonnegative because an $(r,0)$-form in
dimension $r+1$ is decomposable.  If the coefficient curvature
dominates $c\,f^*(i\,dt\wedge d\bar t)$ on a patch where $u(0)$
is nonzero, the right side is strictly positive: its integral on that
patch is bounded below by the positive integral of the squared
relative form times $c\,i\,dt\wedge d\bar t$.
Applying this argument on discs in every base tangent direction proves
horizontal strictness.  When $r=0$ and the full direct image has rank
one, the smooth map has one-point fibers, and its line metric is the
coefficient metric itself.

We now apply this calculation to the theorem.  Write $\Omega_h$ for
the smooth positive volume form on $X$ determined by $h$.  In a
coordinate frame $\partial/\partial x_1\wedge\cdots\wedge
\partial/\partial x_n$ of $-K_X$, it has density equal to that
frame's squared norm.  Its pullback to $Z$ contains the squared
Jacobian factor and vanishes along $E_Z$ to twice its order.

Choose a nonempty Zariski open $U$ over which $g$ is smooth and
base change holds.  After shrinking $U$, the canonical section
$s_{E_Z}$ induces an
isomorphism
\begin{equation}\label{g:volume-adjoint-identification}
 g_*(K_{Z/Y}+\pi^*L)|_U
 =g_*\mathcal O_Z(E_Z)|_U\otimes(-K_Y)|_U
 \simeq(-K_Y)|_U.
\end{equation}
Indeed, the first direct image on the right is generically a line
by \eqref{g:volume-rank-zero}, and $s_{E_Z}$ is a nonzero generic
section of that line.  Local freeness, base change, and its
nonvanishing as a frame hold on a smaller open set.  When
\eqref{g:volume-horizontal-patch} is assumed, the chosen ordinary
patch still meets its inverse image.  We also remove
$\operatorname{Supp}A_0$ from $U$.

In coordinates $y=(y_1,\ldots,y_d)$ on $U$, let $dV_y$ denote the
Euclidean volume with the same normalization used for top forms, and
write
\begin{equation}\label{g:volume-density}
 g_*\pi^*\Omega_h=H(y)\,dV_y.
\end{equation}
By \eqref{g:volume-adjoint-identification}, $H$ is exactly the squared
$L^2$ norm of the corresponding frame of $-K_Y$.  It is smooth and
positive on $U$, and $-\log H$ is plurisubharmonic by smooth
direct-image positivity.  If \eqref{g:volume-horizontal-patch} holds,
the section corresponding to that frame is
nonzero wherever the Jacobian section is nonzero.  Therefore
\eqref{g:volume-curvature-bound}, applied with coefficient
$\mathcal L=\pi^*L$, and
\eqref{g:volume-horizontal-patch} make $-\log H$ strictly
plurisubharmonic at an image point of the patch.  Smoothness then gives
strict positivity on an ordinary neighborhood of that point.

\subsection{Extension and unweighted integrability}

The remaining task is to cross the omitted base locus.  Extension of a
plurisubharmonic weight requires an upper bound; integrability requires
an upper bound on its exponential with the opposite sign.  The
ramification calculation provides the first, while the discrepancy
calculation provides the second.

On a coordinate neighborhood in $Y$, choose local equations $t_Q$ for
the finitely many components of $A_0$.  On its intersection with $U$
define
\begin{equation}\label{g:volume-weight}
 \Phi=-\log H+\sum_Q a_Q\log|t_Q|^2.
\end{equation}
These are the weights of a metric on $-K_Y+A_0$: its squared frame
norm is $H\prod_Q|t_Q|^{-2a_Q}$.  The assertion follows either by
the frame transition rule or after clearing the denominators of
$A_0$.  On $U$ the added terms are pluriharmonic, so the metric has
the curvature already established.

Let $Q$ be a prime component of $Y\setminus U$.  Choose a prime $E$
over $Q$ attaining the minimum in
\eqref{g:volume-correction}.  At general points of $E$ and $Q$, both
are smooth and $E\to Q$ is a submersion.  Choose the resolution to be
an isomorphism there.  After absorbing a nonvanishing unit into a
local coordinate, the map has the form
\[
 (u_1,\ldots,u_n)\longmapsto
 (u_1^{e_E},u_2,\ldots,u_d),
 \qquad Q=\{y_1=0\}.
\]
On a smaller coordinate polydisc the density of $\pi^*\Omega_h$ is
bounded below by a positive constant times $|u_1|^{2a_E}$.
Integrate over one fixed polydisc in the remaining fiber coordinates.
The change of variables $y_1=u_1^{e_E}$, whose real Jacobian is
$e_E^2|u_1|^{2(e_E-1)}$, yields
\begin{equation}\label{g:volume-density-lower-bound}
 H(y)\geq c_Q\,|y_1|^{2b_Q},
 \qquad b_Q=\min_{f(E)=Q}
                  \left(\frac{1+a_E}{e_E}-1\right).
\end{equation}
All other local factors in \eqref{g:volume-weight} are bounded above
and below near the chosen general point.  Since
$b_Q-a_Q=b_Q-\max\{0,b_Q\}\leq0$, the lower bound in
\eqref{g:volume-density-lower-bound} makes $\Phi$ bounded above
there.  The removable-singularity theorem for plurisubharmonic
functions extends it across a general open subset of each omitted
divisor.

The remaining omitted analytic set has codimension at least two in
the smooth base.  The usual Hartogs argument supplies the upper
bound needed there as well.  In a small polydisc, choose one
coordinate direction with circle boundary disjoint from that set.
For transverse parameters outside its projection the entire disc
misses the set, so the submean inequality bounds $\Phi$ by its bound
on the circle.  The exceptional parameter set has measure zero;
applying the submean inequality on a small ball about any remaining
point where $\Phi$ is defined gives the same bound there.  Taking the upper
semicontinuous extension proves removability.  The extensions agree
in overlapping frames by uniqueness, and define a semipositive
metric on $-K_Y+A_0$.

To prove integrability, use the projection formula on the smooth
locus, rather than pulling back the already integrated density:
\begin{equation}\label{g:volume-projection-formula}
 e^{-\Phi}\,dV_y
 =g_*\left(\prod_Q|t_Q\circ g|^{-2a_Q}\,
                      \pi^*\Omega_h\right).
\end{equation}
Take a further resolution $Z'\to W$ on which the Jacobian divisor
and the pullback of $A_0$ have simple normal crossings, and denote
its maps by $\pi':Z'\to X$ and $g':Z'\to Y$.  On a normal-crossing
coordinate chart the upstairs measure on the right of
\eqref{g:volume-projection-formula} is a positive bounded smooth
factor times
\[
 \prod_{j=1}^s |z_j|^{2\gamma_j}\,dV_z,
 \qquad
 \gamma_j=\operatorname{coeff}_{\{z_j=0\}}
                 \bigl(K_{Z'/X}-(g')^*A_0\bigr)>-1,
\]
where the strict inequalities are
Lemma~\ref{g:volume-subklt}.  Each one-variable integral
$\int_0^\varepsilon r^{2\gamma_j+1}\,dr$ is finite, so this measure
is locally integrable on $Z'$.  Properness makes its pushforward
locally finite on $Y$.  Formula~\eqref{g:volume-projection-formula}
then proves $e^{-\Phi}\in L^1$ off the omitted analytic set, with
finite integral on every relatively compact coordinate neighborhood.
The extension changes the weight only on a set of Lebesgue measure
zero and therefore has the same integrability.  This proves
Theorem~\ref{g:toroidal-volume}.

\begin{corollary}\label{g:correction-iitaka-zero}
For the correction in Theorem~\ref{g:toroidal-volume},
$\kappa(Y,A_0)=0$.  More precisely, for every positive integer $m$
for which $mA_0$ is integral,
\[
 H^0(Y,\mathcal O_Y(mA_0))
   =\mathbb C\,s_{mA_0}.
\]
\end{corollary}

\begin{proof}
A section divided by $s_{mA_0}$ is a rational function $q$ on $Y$
with poles bounded by $mA_0$.  The rational function $f^*q$ on
$W$ has only $p$-exceptional poles.  Viewed in
$\mathbb C(W)=\mathbb C(X)$, it consequently has no pole along any
prime of the smooth variety $X$.  Normality extends it regularly
across the remaining subset of codimension at least two.  It is
constant because $X$ is connected and projective.  Dominance of
$f$ makes $q$ constant as well.  The canonical section exists since
$A_0\geq0$, giving the asserted one-dimensional space and Iitaka
dimension.
\end{proof}
